%% !TEX root = manuscript.tex

\chapter*{Introduction}

J'ai réalisé mon alternance de Master 2 MIASHS pour l'année 2025-2026 au sein de la Holding SPA
du groupe Angelotti, un acteur de l'aménagement et de la promotion immobilière implanté en
Occitanie et en région PACA. Cette deuxième année prolonge une première année d'alternance dans
la même entreprise, consacrée à la migration du CRM Vadimm, et ce mémoire en retrace les
principales missions ainsi que la démarche que j'ai adoptée.

Cette année a coïncidé avec l'étape la plus lourde de la modernisation du système
d'information : le remplacement de l'ERP métier, qui passait de Grimmo à la solution SPO éditée
par Salvia. Un changement d'ERP ne se limite pas à changer d'outil, il change aussi le modèle
de données de l'entreprise : les objets ne portent plus les mêmes noms, ne s'articulent plus de
la même façon et n'acceptent plus les mêmes valeurs. Ce chantier m'a conduite d'un rôle
d'assistance opérationnelle vers des missions d'ingénierie, centrées sur la reprise des données
et sur la communication entre les systèmes.

Rattachée à la Holding SPA, j'ai travaillé au contact direct des services métiers, la direction
financière, le service juridique et les équipes commerciales, tout en intervenant sur les
outils du système d'information. Cette position intermédiaire a eu une conséquence pratique
constante : la formulation du besoin fait partie du travail. Un service ne demande presque
jamais un livrable technique, il décrit une gêne, et c'est à moi de la traduire en question
calculable.

Mon activité s'est ainsi répartie sur deux versants complémentaires du métier de la donnée.
D'un côté, un travail d'ingénierie, c'est-à-dire reconstituer dans le nouvel ERP le socle de
données du groupe. De l'autre, un travail d'analyse et de restitution, avec le maintien du
système décisionnel puis la conception d'une application d'aide à la décision pour la direction
financière. Au fil de ces missions, une même difficulté est revenue : avant de pouvoir analyser
une donnée, il fallait presque toujours commencer par la reconstruire. C'est autour de ce
constat que s'articule la problématique de ce mémoire, que je pose au chapitre~1.

Après un premier chapitre de cadrage, je présente les deux grandes missions qui ont rythmé mon
alternance. La Mission~1 (chapitre~2) porte sur la reprise des données vers le nouvel ERP,
d'abord les opérations puis les tiers. La Mission~2 (chapitre~3) porte sur le projet de Data
Science conduit au second semestre, le Copilote Financier. Je présente ensuite les activités
complémentaires menées en parallèle sur les deux semestres (chapitre~4), avant un retour
réflexif sur mon organisation, les compétences acquises, l'adéquation avec le Master MIASHS, le
positionnement RSE de l'entreprise et mon projet professionnel (chapitre~5).
